\chapter{Results \& Discussions}
This chapter presents and discusses the results obtained by running the proposed pipeline on the reference test suite of twelve graphs across six configurations (Cascade baseline plus Nova-Alpha through Nova-Epsilon). It reports the degradation summary, the IPS comparison of Nova-Alpha against Cascade at the same clock, and a detailed root-cause analysis of MosaicNet\_W8A8. The inferences drawn from these results are then discussed.

\section{Contents of this chapter}
All the results obtained for the project objectives are discussed in this chapter. The chapter contains the following sections.
\begin{enumerate}
\item Simulation results - degradation summary across all six configurations.
\item Experimental results - per-graph IPS comparison of Nova-Alpha vs.\ Cascade.
\item Performance Comparison - root-cause analysis of MosaicNet\_W8A8 across Cascade, Nova-Alpha and Nova-Gamma.
\item Inferences drawn from the results obtained.
\end{enumerate}

\subsection{Degradation Summary}
The degradation summary across all six configurations is shown in Table \ref{c5:degradation}. Each option was run on the full test suite of twelve graphs. The Cascade configuration is the baseline against which all Nova variants are compared.

The key observations from the degradation summary are:
\begin{itemize}
\item Nova-Alpha (running at the same clock frequency as Cascade) shows the most degradations - 7 out of 12 tests.
\item Nova-Delta and Nova-Epsilon (1037 MHz) show only 1 degradation each, which is the best result among the Nova variants.
\item Higher clock frequency on the Nova-Delta and Nova-Epsilon configurations compensates for most architectural differences.
\item All regressions passed (zero failures); only performance degradations were observed, not functional failures.
\end{itemize}

% IMAGE PLACEHOLDER: A bar chart of "Degraded count per option" can be added here.

\section{Tables in thesis}
\begin{itemize}
	\item All Table Caption should be in Sentence Case, TNR 10 Pt. It should be of the Format:
	\begin{itemize}
		\item Table 1.1 Results of the experiment ….(Centered)
	\end{itemize}
	\item It should be cited as Table 1.1.
	\item Caption should appear above the Table.
	\item Table Header and the entries should be of Font TNR 10 Pt, Justified.
	\item For wider Table, the page orientation can be Landscape.
	\item For Larger Table, it can run to pages and the header should be repeated for each page of the Table.
	\item Table must be adjusted to fit in the page and no single row is left out for a new page.
\end{itemize}

The degradation summary across all six options is presented in Table \ref{c5:degradation} below.

\begin{table}[htb]
\fontsize{10}{12}\selectfont
\centering
\caption{Degradation summary across configurations (12 tests each)}
\label{c5:degradation}
\begin{tabular}{|p{4cm}|c|c|c|c|}
	\hline
	\textbf{Option (cores)}& \textbf{Total Tests} & \textbf{Pass} & \textbf{Fail} & \textbf{Degraded}\\
	\hline
	\textbf{Cascade (baseline)}   & 12    & 12 &   0 & ---\\\hline
	\textbf{Nova-Alpha}&   12  & 12   & 0 & 7\\\hline
	\textbf{Nova-Beta} & 12 & 12&  0 & 2\\\hline
	\textbf{Nova-Gamma}    &12 & 12&  0 & 3\\\hline
	\textbf{Nova-Delta}&   12  & 12 & 0 & 1\\\hline
	\textbf{Nova-Epsilon}& 12  & 12   & 0 & 1\\
	\hline
\end{tabular}
\end{table}

\subsection{IPS Comparison: Nova-Alpha vs.\ Cascade (Same Clock)}
The per-graph IPS percentage difference of Nova-Alpha relative to Cascade is shown in Table \ref{c5:ipsdiff}. This comparison isolates the architectural effect because both configurations run at the same clock frequency.

\begin{table}[htb]
\fontsize{10}{12}\selectfont
\centering
\caption{Per-graph IPS \% difference of Nova-Alpha vs.\ Cascade (same clock)}
\label{c5:ipsdiff}
\begin{tabular}{|p{6cm}|c|}
	\hline
	\textbf{Graph}& \textbf{IPS \% Diff vs.\ Cascade}\\
	\hline
	LLaMA2\_1B\_Prefill\_AR128 & +9.7\%\\\hline
	LLaMA2\_3B\_Prefill\_AR128 & +6.8\%\\\hline
	MobileVision\_MR\_4x & +2.2\%\\\hline
	LLaMA2\_3B\_Decode\_NativeKV & +0.6\%\\\hline
	MobileVision\_LR\_4x & +0.4\%\\\hline
	MobileBERT\_W8A8 & -1.1\%\\\hline
	VideoDenoise\_HD\_W8A16 & -3.9\%\\\hline
	LLaMA2\_1B\_Decode\_AR1 & -6.1\%\\\hline
	EDSR\_SuperRes\_W8A8 & -6.4\%\\\hline
	MobileDetSSD\_W8A8 & -11.5\%\\\hline
	MobileNetV4\_W8A8 & -16.4\%\\\hline
	MosaicNet\_W8A8 & -24.1\%\\
	\hline
\end{tabular}
\end{table}

% IMAGE PLACEHOLDER: Insert the horizontal bar chart from PPT slide 24 (Nova-Alpha vs Cascade IPS percent difference per graph) here.
\begin{figure}[htb]
\centering
	\includegraphics[scale=1]{Figures/universe}
	\caption{Nova-Alpha vs.\ Cascade -- IPS \% difference per graph}
	\label{fig:ipsdiff}
\end{figure}

\subsection{Root-Cause Analysis: MosaicNet\_W8A8}
MosaicNet\_W8A8 is the worst-degraded test on Nova-Alpha. To isolate the cause, the IPS and the percentage difference of HMX\_STALLS and HVX\_STALLS were compared across Cascade (787 MHz), Nova-Alpha (787 MHz) and Nova-Gamma (1037 MHz). The IPS values observed were 487.7, 370.4 and 459.5 respectively, corresponding to $-24.1\%$ on Nova-Alpha and $-5.8\%$ on Nova-Gamma relative to Cascade. The HVX\_STALLS percentage difference was $-113\%$ on both Nova-Alpha and Nova-Gamma, while HMX\_STALLS percentage difference was approximately $+54\%$ to $+55\%$ (an improvement) on the Nova variants.

% IMAGE PLACEHOLDER: Insert the dual-panel chart from PPT slide 25 (top: IPS bar chart for Cascade / Nova-Alpha / Nova-Gamma; bottom: HMX_STALLS and HVX_STALLS percent difference) here.
\begin{figure}[htb]
\centering
	\includegraphics[scale=1]{Figures/universe}
	\caption{Root-cause analysis of MosaicNet\_W8A8 across Cascade, Nova-Alpha and Nova-Gamma}
	\label{fig:rca}
\end{figure}

\section{Math equation in thesis}
All equations should be written using equation editor or an equivalent tool.
\begin{itemize}
	\item Equations should be numbered as : 1.1, 1.2 ...
	\item Equation should be Centered, 12 Pt, TNR.
	\item Equation number should be right Justified.
	\item It should be cited as Eqn. 1.1.
   \item If the sentence starts by citing an equation, then it should be written as Equation 1.1. For example, Equation 5.1 states the percentage difference computation used by the comparison engine.
\end{itemize}

The percentage difference used to flag degradations and to populate the comparison workbook is given in Eqn.\ \ref{eqn:pctdiff5}.
\begin{equation}
\label{eqn:pctdiff5}
\Delta_{m,o} = \frac{m_o - m_{\text{baseline}}}{m_{\text{baseline}}} \times 100
\end{equation}

The IPS speed-up of Nova-Gamma over Nova-Alpha is given by Eqn.\ \ref{eqn:speedup}, where the higher clock frequency on Nova-Gamma compensates for most of the IPS deficit observed on Nova-Alpha.
\begin{equation}
\label{eqn:speedup}
\text{Speed-up} = \frac{IPS_{\text{Nova-Gamma}}}{IPS_{\text{Nova-Alpha}}}
\end{equation}

\section{Interpretation of results}
The interpretation of the results presented above is summarized as follows.
\begin{itemize}
\item At 787 MHz (Nova-Alpha), HVX stalls explode (\(-113\%\)) -- HVX is severely starved.
\item At 1037 MHz (Nova-Gamma), HVX stalls remain similar in absolute terms but the higher clock compensates, and IPS recovers from \(-24.1\%\) to \(-5.8\%\).
\item HMX stalls actually improve on Nova (fewer HMX stall cycles), confirming that HMX is not the bottleneck on this graph.
\item The remaining \(-5.8\%\) at 1037 MHz is due to HVX stalls that the higher clock cannot fully compensate, pointing engineers towards the HVX scheduling path as the next investigation step.
\item Across the full test suite, higher clock frequency configurations (Nova-Delta and Nova-Epsilon at 1037 MHz) reduce degradations to a single test, confirming that frequency-driven configurations are preferable for the workloads in scope.
\end{itemize}

\vspace{2.5cm}

 \textbf{The chapters should not end with figures, instead bring the paragraph explaining about the figure at the end followed by a summary paragraph.}

In summary, this chapter has presented the regression-level degradation summary across six configurations, the per-graph IPS comparison of Nova-Alpha against Cascade, and a detailed root-cause analysis of the worst-degraded test (MosaicNet\_W8A8). The pipeline successfully reduced the time required to obtain these results from several hours of manual work to approximately two minutes per regression. The next chapter concludes the report and outlines the future scope of the work.
